\begin{center}
    {\Large\bfseries Комбинаторная геометрия}
\end{center}


\problem{1}{Выпуклый $N$-угольник разбит непересекающимися диагоналями на треугольники. Полученные треугольники раскрашены в чёрный и белый цвета так, что любые два треугольника с общей стороной имеют разные цвета. Для каждого $N$ найдите наибольшее возможное значение разности между количеством белых и количеством чёрных треугольников.}

% Источник: Турнир городов, 2002/03, осенний тур, основной вариант,
% 10--11 классы, задача 3.
% https://problems.ru/view_problem_details_new.php?id=98594

\medskip

\problem{2}{На плоскости расположено
\[
    \left\lfloor\frac{4n}{3}\right\rfloor
\]
прямоугольников, стороны которых параллельны осям координат. Известно, что каждый прямоугольник пересекается хотя бы с $n$ другими. Докажите, что найдётся прямоугольник, пересекающийся со всеми остальными.}

% Источник: региональный этап ВсОШ, 2001/02,
% 9 класс, задача 4; автор В.\,Л.~Дольников.
% https://problems.ru/view_problem_details_new.php?id=110102

\medskip

\problem{3}{На окружности отмечено $2N$ точек. Известно, что через любую точку внутри окружности проходит не более двух хорд с концами в отмеченных точках.

Паросочетанием назовём набор из $N$ хорд, в котором каждая отмеченная точка является концом ровно одной хорды. Паросочетание назовём чётным, если количество точек пересечения его хорд чётно, и нечётным в противном случае. Найдите разность между количеством чётных и количеством нечётных паросочетаний.}

% Источник: региональный этап ВсОШ, 2011/12,
% 10 класс, задача 4; автор В.~Шмаров.
% https://problems.ru/view_problem_details_new.php?id=116590

\medskip

\problem{4}{На плоскости отмечено $2n$ точек, никакие три из которых не лежат на одной прямой. Половина точек окрашена в красный цвет, а половина~--- в синий. Прямую назовём уравновешивающей, если она проходит через одну красную и одну синюю точку, а в каждой из двух открытых полуплоскостей, на которые она делит плоскость, находится поровну красных и синих точек. Докажите, что существуют по крайней мере две уравновешивающие прямые.}

% Источник: USA Mathematical Olympiad, 2005, задача 2.
% https://artofproblemsolving.com/wiki/index.php/2005_USAMO_Problems/Problem_2


\problem{5}{На плоскости отмечено несколько точек, каждая из которых окрашена в синий, жёлтый или зелёный цвет. На любом отрезке, соединяющем две одноцветные точки, нет других точек этого цвета, но есть хотя бы одна отмеченная точка другого цвета. Каково наибольшее возможное количество отмеченных точек?}

% Источник: заключительный этап ВсОШ, 2006/07,
% 11 класс, задача 6; автор М.\,В.~Мурашкин.
% https://problems.ru/view_problem_details_new.php?id=111767

\medskip

\problem{6}{Пусть $S$~--- конечное множество, состоящее не менее чем из двух точек плоскости, причём никакие три точки множества $S$ не лежат на одной прямой.

Назовём \emph{мельницей} следующий процесс. Сначала выбираются точка $P\in S$ и прямая $\ell$, проходящая через $P$ и ни через какую другую точку множества $S$. Прямая $\ell$ вращается по часовой стрелке вокруг точки $P$ до тех пор, пока впервые не встретит другую точку $Q\in S$. После этого точка $Q$ становится новым центром вращения, и прямая продолжает вращаться по часовой стрелке вокруг $Q$ до встречи со следующей точкой множества $S$. Процесс продолжается неограниченно.

Докажите, что начальные точку $P$ и прямую $\ell$ можно выбрать так, чтобы каждая точка множества $S$ становилась центром вращения бесконечно много раз.}

% Источник: Международная математическая олимпиада, 2011, задача 2.
% https://www.imo-official.org/problems/2011/eng.pdf

\medskip

\problem{7}{На плоскости дано конечное семейство равных квадратов с соответственно параллельными сторонами. Известно, что среди любых $k+1$ квадратов найдутся два пересекающихся. Докажите, что всё семейство можно разбить не более чем на $2k-1$ непустых подсемейств так, чтобы во всяком подсемействе все квадраты имели общую точку.}

% Источник: заключительный этап ВсОШ, 1994/95,
% 11 класс, задача 4; автор В.\,Л.~Дольников.
% https://problems.ru/view_problem_details_new.php?id=109859

\medskip

\problem{8}{Конфигурацию из $4027$ точек плоскости назовём \emph{колумбийской}, если она состоит из $2013$ красных и $2014$ синих точек, причём никакие три точки конфигурации не лежат на одной прямой.

Проведя несколько прямых, плоскость разбивают на области. Набор проведённых прямых назовём \emph{хорошим} для данной колумбийской конфигурации, если выполнены условия:
\begin{itemize}
    \item ни одна проведённая прямая не проходит через точку конфигурации;
    \item ни одна из полученных областей не содержит точек разных цветов.
\end{itemize}
Найдите наименьшее число $k$, для которого у любой колумбийской конфигурации существует хороший набор из $k$ прямых.}

% Источник: Международная математическая олимпиада, 2013, задача 2.
% https://www.imo-official.org/problems/2013/eng.pdf